\subsection{Actual polynomial logarithms are subharmonic}\label{subsec:fc-line-polynomial-subharmonic}

\begin{lemma}[Polynomial logarithm minus an actual harmonic function]\label{lem:fc-polynomial-log-local-maximum}
\lean{EnvelopeProperty.Berkovich.RelativeSpectrum.FCStandardAffinoid.logPolynomial_sub_harmonic_eventually_eq_of_isLocalMax}
\leanok
\uses{analytic:actual-original-unit-max-all-point-local-constancy, def:fc-actual-harmonic}
For a nonzero polynomial $P$ and an actual harmonic expression
$H=c+\sum_j\gamma_j\log|T-a_j|$, every local maximum of
$\log|P|-H$ on the common nonvanishing locus of $P$ and the $T-a_j$
has an actual open neighbourhood on which this same function is constant.
\end{lemma}

\begin{proof}
Use the original finite rational-unit family consisting of $P$ and the
$T-a_j$, with denominator $1$ for every member, coefficients $1$ and
$-\gamma_j$, and constant $-c$. On the common unit locus its actual affine
value is exactly $\log|P|-H$. Apply the all-point original rational-unit
local maximum theorem to the singleton maximum of this family.
\end{proof}

\begin{theorem}[Actual polynomial logarithms are subharmonic]\label{thm:fc-polynomial-log-subharmonic}
\lean{EnvelopeProperty.Berkovich.RelativeSpectrum.isLineSubharmonicOn_logAbs_polynomial}
\leanok
\uses{def:fc-actual-line-subharmonic, lem:fc-line-boundary-comparison,
  lem:fc-polynomial-log-local-maximum}
For every polynomial $P$, the actual function $x\mapsto\log|P(x)|$, with
$\log0=-\infty$, is subharmonic on every domain in the analytic line.
\end{theorem}

\begin{proof}
The actual extended logarithmic seminorm evaluation is continuous and
never $+\infty$. The zero polynomial gives the constant $-\infty$ function.
For $P\ne0$, compare with an actual harmonic $h$ on an actual standard
affinoid $U$. The difference $F=\log|P|-h$ is continuous as an extended-real
function on $U$, and is nonpositive on its actual boundary by hypothesis.
At a positive interior local maximum, $P$ cannot vanish; the hole
coordinates are positive everywhere on $U$. The preceding actual unit
calculation therefore supplies the required local plateau for $F$.
Compact connected boundary comparison gives $F\le0$ on $U$, which is the
literal harmonic comparison inequality. No analytic maximum-principle
certificate is assumed.
\end{proof}
